\section{Related Work}
\label{sec:related}

\paragraph{Instruction tuning.} Alpaca \cite{alpaca} popularized the
recipe of fine-tuning an open base model on a small synthetic instruction
dataset generated by a stronger teacher, and the Self-Instruct line
\cite{selfinstruct} established that LM-generated instruction data is good
enough to train a useful assistant. UltraChat and follow-on dialog datasets
extended this to multi-turn and long-context settings. Our pipeline is
firmly in this tradition: a stronger teacher model (Claude Sonnet 4.6)
distills source material into instruction pairs that an open student
(Qwen~2.5\,7B-Instruct) learns from. The difference is target: we tune for
\emph{register} (Summa-style quaestio, Augustinian address, CCC citation
form), not for general helpfulness.

\paragraph{Parameter-efficient fine-tuning.} LoRA \cite{lora} factors weight
updates into low-rank adapters, drastically reducing the number of
trainable parameters and the memory required to fine-tune large models;
DoRA decomposes the update into magnitude and direction for additional
sample efficiency. We adopt LoRA on the top 16 transformer layers under
MLX, which together with 8-bit weight quantization brings end-to-end
training of a 7B model into the 12--13\,GB memory budget of a 48\,GB
unified-memory laptop. Updating only 2.6M of 7.6B parameters appears
sufficient for stylistic transfer in our setting.

\paragraph{Preference optimization.} DPO \cite{rafailov2023dpo} reframes
RLHF \cite{ouyang2022instructgpt} as a single likelihood objective on
preference pairs, sidestepping a separate reward model. ORPO \cite{orpo}
folds reference-free preference learning into a single training stage with
an SFT-like loss. GRPO \cite{grpo} (introduced in DeepSeekMath) generalizes
to grouped relative-preference signals. We tried DPO and ORPO before SFT
in this project, on a smaller (1.5B, 4-bit) base, and observed the
asymmetric-reward failure mode: preference-objective validation metrics
looked excellent (DPO val accuracy~1.0, margin~7.5 nats) while generations
were nearly byte-identical to the base for most prompts.
Section~\ref{sec:discussion} returns to why we believe SFT-first on a
larger Q8 base broke this pattern.

\paragraph{On-device LLM fine-tuning.} MLX \cite{mlx} provides a
unified-memory-aware NumPy-like framework for Apple Silicon, with native
support for LoRA fine-tuning of quantized weights. Recent reports show
that 7B-class models can be supervised fine-tuned on 48--64\,GB M-series
machines in minutes rather than hours, removing a key barrier to small
research projects. Our work is one such report.

\paragraph{Stylistic and domain adaptation.} A long line of work
demonstrates that fine-tuning shifts surface style as readily as it shifts
factual content, often more readily. Stylistic transfer in LLMs has been
explored for legal English, medical English, and literary translation. We
focus on a niche register (thirteenth-century scholastic and
fourth/fifth-century patristic theology in English translation) and
adopt teacher distillation precisely because hand-curated data at the
register level is impractical at small scale.

\paragraph{Religious-text NLP.} The bulk of corpus work on religious texts
sits in parallel-translation studies (Bible translations, Quran translations)
and in lexicon-and-grammar tasks for low-resource liturgical languages.
LLM-era work on the Catechism, Summa, or Augustinian corpus appears to be
sparse; we are not aware of prior work explicitly fine-tuning a
contemporary chat model to adopt the Summa's \emph{quaestio} surface form
together with CCC paragraph grounding. We note this as a gap rather than
as a claim of novelty against an exhaustive search.
